% part: intuitionistic-logic
% chapter: soundness-completeness
% section: lindenbaum

\documentclass[../../../include/open-logic-section]{subfiles}

\begin{document}

\olfileid[pt]{int}{sc}{lin}

\olsection{Lema de Lindenbaum}

O teorema da completude para a lógica intuicionista é provado supondo
$\Gamma \Proves/ !A$ e construindo um modelo tal que
$\mSat{M}{\Gamma}$ e~$\mSat/{M}{!A}$.

Na lógica clássica, a relação de !!{derivability} pode ser reduzida à
noção de consistência, pois !!a{formula}~$!A$ é !!{derivable} a partir
de um conjunto de !!{formula}s sse esse conjunto, juntamente com a
negação de~$!A$, é inconsistente. Isso não é possível na lógica
intuicionista. Na lógica intuicionista, se $\lnot!A$ é inconsistente,
obtemos apenas $\Proves \lnot\lnot !A$. Como
$\lnot\lnot!A \lif !A$ não vale intuicionisticamente em geral, não
podemos concluir~$\Proves !A$.

Assim, ao construir o modelo~$\mModel{M}$, precisaremos acompanhar a
não !!{derivability} da !!{formula}~$!A$ e, portanto, não poderemos
usar um conjunto completo $\Gamma^* \supseteq \Gamma$ para construir o
modelo $\mModel{M}$, pois, em todo conjunto completo $\Gamma^*$, temos
$\Gamma^* \Proves !A \lor \lnot !A$.

Em vez de usar um conjunto completo~$\Gamma^*$, empregaremos a noção
de conjunto primo de fórmulas:

\begin{defn}\ollabel{defn:prime}
  Um conjunto de !!{formula}s~$\Gamma$ é \emph{primo} sse
  \begin{enumerate}
  \item\ollabel{defn:prime1} $\Gamma$ é consistente, O.s.\@,
    $\Gamma \Proves/ \lfalse$;
  \item\ollabel{defn:prime2} se $\Gamma \Proves !A$, então
    $!A \in \Gamma$; e
  \item\ollabel{defn:prime3} se $!A \lor !B \in \Gamma$, então
    $!A \in \Gamma$ ou $!B \in \Gamma$.
  \end{enumerate}
\end{defn}

\begin{lem}[Lema de Lindenbaum]
  \ollabel{lem:lindenbaum} Se $\Gamma \Proves/ !A$, existe
  $\Gamma^* \supseteq \Gamma$ tal que $\Gamma^*$ é primo e
  $\Gamma^* \Proves/ !A$.
\end{lem}

\begin{proof}
  Seja $!B_1 \lor !C_1$, $!B_2 \lor !C_2$, \dots, uma enumeração de
  todas as !!{formula}s da forma~$!B \lor !C$. Definiremos uma
  sequência crescente de conjuntos de !!{formula}s~$\Gamma_n$, na
  qual cada $\Gamma_{n+1}$ é definido como $\Gamma_n$ juntamente com
  uma nova !!{formula}. $\Gamma^*$ será a união de todos os
  $\Gamma_n$. As novas !!{formula}s são escolhidas de modo a garantir
  que $\Gamma^*$ seja primo e que ainda $\Gamma^* \Proves/ !A$. Isso
  significa que, em cada etapa, devemos encontrar a primeira
  disjunção~$!B_i \lor !C_i$ tal que:
  \begin{enumerate}
  \item\ollabel{gamma-1} $\Gamma_n \Proves !B_i \lor !C_i$
  \item\ollabel{gamma-2} $!B_i \notin \Gamma_n$ e $!C_i \notin \Gamma_n$
  \end{enumerate}
  Acrescentamos a $\Gamma_n$ ou $!B_i$, se
  $\Gamma_n \cup \{!B_i\} \Proves/ !A$, ou $!C_i$, caso contrário.
  Teremos de mostrar que isso funciona. Por ora, definamos $i(n)$ como
  o menor $i$ tal que \olref{gamma-1} e~\olref{gamma-2} valham.
  
  Define $\Gamma_0 = \Gamma$ and
  \[
  \Gamma_{n+1} = \begin{cases}
    \Gamma_n \cup \{!B_{i(n)}\} &
    \text{se $\Gamma_n \cup \{!B_{i(n)}\} \Proves/ !A$} \\
    \Gamma_n \cup \{!C_{i(n)}\} & \text{caso contrário}
    \end{cases}
  \]
  Se $i(n)$ não está definido, O.s.\@, se, sempre que
  $\Gamma_n \Proves !B \lor !C$, ou $!B \in \Gamma_n$ ou
  $!C \in \Gamma_n$, estabelecemos $\Gamma_{n+1} = \Gamma_n$.
  Agora, seja $\Gamma^* = \bigcup_{n=0}^\infty \Gamma_n$.

  Primeiro mostramos que, para todo $n$, $\Gamma_n \Proves/ !A$.
  Procedemos por indução sobre~$n$. Para $n = 0$, a afirmação vale pela
  hipótese do teorema, O.s.\@, $\Gamma \Proves/ !A$. Se $n>0$,
  precisamos mostrar que, se $\Gamma_n \Proves/ !A$, então
  $\Gamma_{n+1} \Proves/ !A$. Se $i(n)$ não está definido,
  $\Gamma_{n+1} = \Gamma_n$ e nada há a provar. Suponha, então, que
  $i(n)$ esteja definido. Por simplicidade, seja $i = i(n)$.
  
  Provaremos a contrapositiva da afirmação. Suponha
  $\Gamma_{n+1} \Proves !A$. Por construção,
  $\Gamma_{n+1} = \Gamma_n \cup \{!B_i\}$ se
  $\Gamma_n \cup \{!B_i\} \Proves/ !A$; caso contrário,
  $\Gamma_{n+1} = \Gamma_n \cup \{!C_i\}$. Evidentemente, não pode ser
  o primeiro caso, pois então $\Gamma_{n+1} \Proves/ !A$. Logo,
  $\Gamma_n \cup \{!B_i\} \Proves !A$ e
  $\Gamma_{n+1} = \Gamma_n \cup \{!C_i\}$. Pela definição de $i(n)$,
  temos $\Gamma _n \Proves !B_i \lor !C_i$. Temos também
  $\Gamma_n \cup \{!B_i\} \Proves !A$ e
  $\Gamma_{n+1} = \Gamma_n \cup \{!C_i\} \Proves !A$. Portanto,
  $\Gamma_n \Proves !A$, que é o que queríamos mostrar.

  Se $\Gamma^* \Proves !A$, existiria algum subconjunto finito
  $\Gamma' \subseteq \Gamma^*$ tal que $\Gamma' \Proves !A$. Cada
  $!D \in \Gamma'$ deve pertencer a $\Gamma_i$ para algum~$i$. Seja
  $n$ o maior desses índices. Como $\Gamma_i \subseteq \Gamma_n$ se
  $i \le n$, temos $\Gamma' \subseteq \Gamma_n$. Mas então
  $\Gamma_n \Proves !A$, contrariando nossa prova acima de que
  $\Gamma_n \Proves/ !A$.

  Por fim, mostramos que $\Gamma^*$ é primo, O.s.\@, satisfaz as
  condições \olref{defn:prime1}, \olref{defn:prime2}
  e~\olref{defn:prime3} de \olref{defn:prime}.

  Primeiro, $\Gamma^* \Proves/ !A$; portanto, $\Gamma^*$ é consistente
  e \olref{defn:prime1} vale.

  Mostramos agora que, se $\Gamma^* \Proves !B \lor !C$, então ou
  $!B \in \Gamma^*$ ou $!C \in \Gamma^*$. Isso prova
  \olref{defn:prime3}, pois, se $!B \lor !C \in \Gamma^*$, então
  também $\Gamma^* \Proves !B \lor !C$. Suponha, portanto,
  $\Gamma^* \Proves !B \lor !C$, mas $!B \notin \Gamma^*$ e
  $!C \notin \Gamma^*$. Como $\Gamma^* \Proves !B \lor !C$,
  $\Gamma_n \Proves !B \lor !C$ para algum~$n$. $!B \lor !C$ aparece
  na enumeração de todas as disjunções, digamos como
  $!B_j \lor !C_j$. $!B_j \lor !C_j$ satisfaz as propriedades da
  definição de $i(n)$: temos $\Gamma_n \Proves !B_j \lor !C_j$,
  enquanto $!B_j \notin \Gamma_n$ e $!C_j \notin \Gamma_n$. Em cada
  etapa, ao menos uma disjunção $!B_i \lor !C_i$ a menos satisfaz as
  condições (pois, em cada etapa, acrescentamos ou $!B_i$ ou $!C_i$);
  assim, em alguma etapa~$m$, teremos $j = i(m)$. Mas então ou
  $!B \in \Gamma_{m+1}$ ou $!C \in \Gamma_{m+1}$, contrariando a
  hipótese de que $!B \notin \Gamma^*$ e $!C \notin \Gamma^*$.

  Suponha agora $\Gamma^* \Proves !B$. Então
  $\Gamma^* \Proves !B \lor !B$. Mas acabamos de provar que, se
  $\Gamma^* \Proves !B \lor !B$, então $!B \in \Gamma^*$. Logo,
  $\Gamma^*$ satisfaz~\olref{defn:prime2} de~\olref{defn:prime}.
\end{proof}

\begin{prob}
Mostre que, se $\Gamma \Proves/ \bot$, então $\Gamma$ é consistente na
lógica clássica, O.s.\@, existe uma valoração que torna verdadeiras
todas as fórmulas de~$\Gamma$.
\end{prob}

\end{document}
